% !TeX root = ../../main.tex
\subsection{Inverting the sine transform and transferring coefficients}
\label{subsec:s08:04}

We must recover the raw potential, including its constant in $x$.
The inverse sine operator produces the dominant coefficient; the
constant of its transformed primitive is smaller, but its size does
not follow from the differential equation alone. We first give an
integral form of the inverse operator, then control its whole path
before extracting coefficients.

For $g\geq h\geq1$ put
\begin{equation}\label{eq:s08:degree}
 D_{g,h}=g_2+h=2g+h-3.
\end{equation}
Define the Borel series of the difference of potentials by
\begin{equation}\label{eq:s08:inverse-borel-definition}
 \mathcal F_h^+(\theta;x,y)
 =\sum_{g\geq h}(2\pi)^{2g}
   (\mathcal F_{g,h}-\mathcal G_{g,h})(\theta;x)
                       \frac{y^{D_{g,h}}}{D_{g,h}!}.
\end{equation}
Its superscript refers to the strictly positive powers of the inverse
operator, not to a positivity assertion. Its parity is $(-1)^{h-1}$.

\begin{lmm}[The inverse sine kernel]\label{lmm:s08:inverse-kernel}
For each fixed $h\geq1$, one has the exact identity
\begin{equation}\label{eq:s08:inverse-kernel}
 \begin{aligned}
 \mathcal F_h^+(\theta;x,y)
 =\sum_{m\geq1}\frac{(-1)^{m-1}}{m^2}\int_0^y
 \bigl[&\partial_x\mathcal L_h(\theta;x+(y-u)/m,u)\\
       &+\partial_x\mathcal L_h(\theta;x-(y-u)/m,u)\bigr]\,\dd u.
 \end{aligned}
\end{equation}
Here $u$ runs along the segment from $0$ to $y$, and the derivatives
are taken before evaluation at the shifted centers. The identity
holds first as a formal series and then by normal convergence on
the common analytic regions near $x=0$.
\end{lmm}
\begin{prf}
The evaluation of $\beta_k$ in Theorem~\ref{thm:s06:lambda-g} gives
$(2\pi)^{2k}\beta_k
=2\sum_{m\geq1}(-1)^{m-1}m^{-2k}$ for $k\geq1$.
In the two shifts in \eqref{eq:s08:inverse-kernel}, only even Taylor
degrees remain. Write that degree as $a=2k-2$ and the input Borel
degree as $N'=2g'+h-2$. The elementary integral
\[
 \int_0^y u^{N'}(y-u)^a\,\dd u
       =\frac{N'!a!}{(N'+a+1)!}y^{N'+a+1}
\]
changes the factorial to $(N'+2k-1)!=D_{g'+k,h}!$.
The two shifts contribute $2m^{-2k}$, and their derivative is
$\partial_x^{2k-1}L_{g',h}=\partial_x^{2k}\mathcal G_{g',h}$.
Thus every coefficient is exactly the $k\geq1$ part of
\eqref{eq:s07:inverse-miura}.

For the analytic assertion, the shifted radius budget is the same
one as in Lemma~\ref{lmm:b:cone-transport}. On compact subsets the
derivatives are bounded uniformly in $m$, and the outer $m^{-2}$
is summable. The regular patches and endpoint bounds of
Proposition~\ref{prp:b:common-domains} continue this normally
convergent identity to the stated regions.
\end{prf}

The positive $m=1$ branch is singular along the entire path when
$y$ reaches $R_\theta-x$. Its short initial segment approaches the
critical center itself. It is therefore essential to treat that
segment separately, rather than insert a fixed-center expansion
under the integral without a uniform bound.

\begin{prp}[The recovered singularity]\label{prp:s08:recovered-singularity}
Fix $h\geq1$ and $1/2<\nu<1$. Uniformly for $x$ in a sufficiently
small complex neighborhood of $0$ and small real $\theta\geq0$,
with $S=R_\theta-x$, one has
\begin{equation}\label{eq:s08:recovered-singularity}
 \begin{aligned}
 \mathcal F_h^+(\theta;x,y)
 ={}&\frac{1}{2}\sqrt{\frac{2}{-w''(x_{*,\theta})}}
       \frac{\displaystyle\int_x^{R_\theta}P_h(\theta,s)\,\dd s}
            {\sqrt{S-y}}\\
 &\quad+O_{h,\nu}(|1-y/S|^{\nu-1})
                         \qquad(y\longrightarrow S).
 \end{aligned}
\end{equation}
The expansion holds on a common dented complex neighborhood, with
an additive error even when the displayed integral vanishes.
The expansion at $-S$ is fixed by parity.
\end{prp}
\begin{prf}
Consider the positive $m=1$ term of
\eqref{eq:s08:inverse-kernel}, with $u=\tau y$ and $0\leq\tau\leq1$.
At its center $s=x+y-u$, the center distance is $S-y+u$
and the endpoint distance remains $S-y$. Use a fixed angular
reduction of the common dented neighborhood, so that $(S,y)$ lies
strictly inside the larger two-distance cone used for the estimates
in Proposition~\ref{prp:b:common-domains}.
Lemma~\ref{lmm:b:cone-transport} preserves this angular margin
under the shift $(x,y)\mapsto(x+(1-\tau)y,\tau y)$.
All constants below are uniform on this smaller complex region.

\smallskip\noindent
\emph{The initial segment.}
Choose a sufficiently large fixed $C>0$ and split the segment at
$|u|=C|S-y|$. On its initial part $\tau<1/2$ once the endpoint
neighborhood is small enough. The three distances satisfy
\[
 \begin{aligned}
 R_\theta-s&=(1-\tau)(S-y)+\tau S,\\
 (R_\theta-s)-u&=S-y,\\
 (R_\theta-s)+u&=(1-\tau)(S-y)+\tau(S+y).
 \end{aligned}
\]
Positive-cone addition bounds the moduli of the first and third
expressions below by a fixed multiple of $|S-y|$, because
$1-\tau\geq1/2$. Their upper bounds follow from
$|u|\leq C|S-y|$. Consequently
\[
 |R_\theta-s|\asymp |S-y|,\qquad
 |(R_\theta-s)\pm u|\asymp |S-y|.
\]
Both normalized endpoint distances $1\pm u/(R_\theta-s)$
are therefore bounded above and away from zero.
As $y\to S$, this is a small critical-center cone with fixed
angular margins. Proposition~\ref{prp:b:common-domains} includes
its holomorphic extension and its patching to the common complex
region. Apply \eqref{eq:b:common-bounds} there with one center
derivative and no Borel derivative:
\[
 |\partial_s\mathcal L_h(\theta;s,u)|
       \leq C_{h,C}|S-y|^{(h-3)/2}.
\]
Multiplication by the path length $O(C|S-y|)$ gives the uniform
initial integral $O_h(|S-y|^{(h-1)/2})$, bounded for all $h\geq1$.
On the remaining segment, positive-cone addition gives
$|S-y+u|\asymp |S-y|+|u|$, so that
$|(S-y)/(S-y+u)|$ is small enough for
the endpoint expansion \eqref{eq:s08:leading-expansion} to apply.

\smallskip\noindent
\emph{The principal term on the remaining segment.}
Differentiate that expansion in its center while holding its Borel
argument fixed. The square root contributes
\[
 \frac12\sqrt{\frac{2}{-w''(x_{*,\theta})}}
               \frac{P_h(\theta,s)}{\sqrt{S-y}}.
\]
Its integral is the term in \eqref{eq:s08:recovered-singularity}.
Indeed, extending the center path to the full segment from $x$ to
$R_\theta$ changes the numerator by
$O_h(|S-y|^{h/2})$, using
$P_h(\theta,s)=O_h(|R_\theta-s|^{(h-2)/2})$.
After division by $\sqrt{S-y}$ this change is bounded.

The derivative of the endpoint value integrates to a difference
of endpoint values and is bounded by their scale $(h-1)/2$.
The derivative of $P_h$ has an integral equal to a difference of
$P_h$ values, multiplied by $\sqrt{S-y}$, and is bounded as well.
For the remaining error, the mixed Cauchy bounds following
\eqref{eq:s08:leading-bounds} give
\[
 |\partial_sE_h(\theta;s,u)|
 \leq C_{h,\nu}|S-y|^{\nu-1}
                    |S-y+u|^{(h-1)/2-\nu}.
\]
The last exponent is greater than $-1$. Its integral is therefore
$O_{h,\nu}(|S-y|^{\nu-1})$, as required.

\smallskip\noindent
\emph{The other branches.}
For the negative $m=1$ branch, and for either branch with $m\geq2$,
the critical endpoint is met only at $u=y$. Near it the first
center derivative has a bound by a constant times
$(|S-y|+|y-u|)^{-1/2}$; the remainder terms are integrable too.
For $m\geq2$ the positive branch has the uniform distance factor
$1-1/m\geq1/2$. Their sums are bounded by the outer $m^{-2}$.
Away from this final segment all shifted arguments range in
regular compact inverse patches, where
Proposition~\ref{prp:b:tube} applies. For negative centers bounded
away from zero, Lemma~\ref{lmm:b:negative-disks} supplies an inverse
radius strictly larger than $R$; compact perturbation preserves the
required margin. Thus those centers use regular patches, not a
positive-center endpoint expansion. This proves the bound on the whole path and its normal
continuation. Parity follows from the defining Borel degrees.
\end{prf}

We use the following two-point form of singularity transfer. Its
short proof also records the error strength actually available here.

\begin{lmm}[Two singularities]\label{lmm:s08:transfer}
Let $1/2<\nu<1$. Suppose a function $H(Y)$ is holomorphic in
$|Y|<1$ and in a doubly dented enlargement as in
\eqref{eq:s08:dented-domains}, with $S=1$. Suppose, at its two
endpoints, respectively,
\[
 H(Y)=A_+(1-Y)^{-1/2}+O(|1-Y|^{\nu-1}),\qquad
 H(Y)=A_-(1+Y)^{-1/2}+O(|1+Y|^{\nu-1}),
\]
where $A_+,A_-\in\mathbb C$. As integers $N\to\infty$,
\begin{equation}\label{eq:s08:two-point-transfer}
 [Y^N]H(Y)=\frac{A_++(-1)^NA_-}{\sqrt{\pi N}}
                                     +O(N^{-\nu}).
\end{equation}
The assertion is uniform for common domains, bounded amplitudes
and uniform error and compact-subset bounds.
\end{lmm}
\begin{prf}
Subtract both displayed inverse square roots. The resulting function
has order $|1\mp Y|^{\nu-1}$ at each endpoint. Deform its Cauchy
contour to two indented contours with tip radii $1/N$. Each tip
contributes $O(N^{-\nu})$. On each outgoing ray take
$1-Y=t e^{\pm i\varphi}$, with a fixed
$\pi/2<\varphi<\pi-\gamma$, and reflect these rays at the
negative endpoint. For small fixed $t>0$ one has
$|Y|^{-N}\leq e^{-cNt}$, $c>0$. Their contributions are bounded by
$C\int_{1/N}^{r}t^{\nu-1}e^{-cNt}\,\dd t=O(N^{-\nu})$.
The remaining contour is outside the unit circle by a fixed margin
and gives an exponentially small term. Finally
$[Y^N](1-Y)^{-1/2}=\binom{2N}{N}/4^N
=(\pi N)^{-1/2}(1+O(N^{-1}))$.
This proves \eqref{eq:s08:two-point-transfer} with the stated
uniformity.
\end{prf}

\begin{rmk}[Transfer method]
This is the contour method of Flajolet--Odlyzko
\cite[\S 2]{FO90}, applied at two isolated singularities. The common
dented regions above supply its complex-domain hypothesis; real-axis
asymptotics alone would not supply it.
\end{rmk}

For $g\geq2$ and $0\leq h\leq g$, write
$T_{g,h}=T_{g,h,0}$, with the integral defined in
\eqref{eq:s07:kappa-integrals}. For $h=0$ we also use
$D_{g,0}=g_2$, extending \eqref{eq:s08:degree} to this case.

\begin{prp}[Recovered intersections]\label{prp:s08:intersections}
For each fixed integer $h\geq0$, as $g\to\infty$,
\begin{equation}\label{eq:s08:intersection-asymptotic}
 T_{g,h}\sim
 \frac{\beta_g D_{g,h}!R^{-D_{g,h}}}
      {J_1(z_{0,1})\sqrt{2\pi D_{g,h}}}
                       \frac{\mathcal B(0)^h}{h!}.
\end{equation}
In particular,
\begin{equation}\label{eq:s08:intersection-ratio}
 \frac{T_{g,h}}{T_{g,0}}
 \sim\frac{D_{g,h}!}{g_2!}R^{-h}
             \frac{\mathcal B(0)^h}{h!}
             \sqrt{\frac{g_2}{D_{g,h}}}.
\end{equation}
For every fixed $h\geq0$ and $n\in\{0,1,2\}$ there is $C_{h,n}>0$
such that, whenever $g\geq\max\{1,h\}$ and $2g-2+n>0$,
with $M=2g-3+h+n$,
\begin{equation}\label{eq:s08:marked-bound}
 |T_{g,h,n}|\leq C_{h,n}\beta_g
                          \frac{M!R^{-M}}{\sqrt{M+1}}.
\end{equation}
\end{prp}
\begin{prf}
For $h\geq1$, set $\theta=x=0$ in
\eqref{eq:s08:recovered-singularity} and use
\eqref{eq:s08:ordered-integrals}. The amplitude multiplying
$(1-y/R)^{-1/2}$ is
\[
 \frac{z_{0,1}}{2\sqrt2 R}\,
                         \frac{\mathcal B(0)^h}{h!}.
\]
At the negative endpoint it is multiplied by $(-1)^{h-1}$.
Since $D_{g,h}$ has parity $h-1$, both endpoints give the same
coefficient: their sum has a factor two. Applying
Lemma~\ref{lmm:s08:transfer} and then
\eqref{eq:s06:beta-asymptotics} gives the right-hand side of
\eqref{eq:s08:intersection-asymptotic} for
$(\mathcal F_{g,h}-\mathcal G_{g,h})(0;0)$.
Here $R=z_{0,1}J_1(z_{0,1})/2$ converts the constant.

\smallskip\noindent
\emph{The primitive constant.}
Proposition~\ref{prp:b:primitive-bound}, with $N=D_{g,h}+1$, gives
\[
 |\mathcal G_{g,h}(0;0)|
 \leq C_h(2\pi)^{-2g}(D_{g,h}+1)!
                  R^{-D_{g,h}-1}(D_{g,h}+1)^{-5/2}.
\]
It is smaller by $O_h(g^{-1})$ than the scale just found.
Its proof uses the geometric primitive on the negative real axis,
not an arbitrary choice of integration constant. The genus-one
logarithm is excluded from that bound and does not affect this
large-genus assertion. Equation~\eqref{eq:s07:specialization}
now identifies $\mathcal F_{g,h}(0;0)$ with $T_{g,h}$.
For $h=0$, the same asymptotic follows directly from
\eqref{eq:s06:exact-vertex} and \eqref{eq:s06:additive-saddle}
with amplitude one. Taking their ratio proves
\eqref{eq:s08:intersection-ratio}.

\smallskip\noindent
\emph{Fixed markings.}
The proof of \eqref{eq:s08:recovered-singularity} is uniform for
complex $x$ near zero. Lemma~\ref{lmm:s08:transfer}, normalized by
$Y=y/(R-x)$, therefore bounds the coefficients of
$\mathcal F_h^+$ by
\[
 C_h(2\pi)^{-2g}D_{g,h}!|R-x|^{-D_{g,h}}
                                      (D_{g,h}+1)^{-1/2}.
\]
Cauchy's formula on $|x|=c/(D_{g,h}+1)$, with a small fixed
$c>0$, costs at most $C_n(D_{g,h}+1)^n$ for $n$ derivatives.
The factor $|R-x|^{-D_{g,h}}$ is at most a constant times
$R^{-D_{g,h}}$ on this circle. For the transformed part and
$n\geq1$, use $\partial_x\mathcal G_{g,h}=L_{g,h}$ and
Lemma~\ref{lmm:b:coefficient-bound} on the same parameter circles.
This gives a bound by
\[
 C_{h,n}(2\pi)^{-2g}(D_{g,h}+1)!R^{-D_{g,h}-1}
                  (D_{g,h}+1)^{-3/2+n-1},
\]
which is smaller than the required scale. For $n=0$ use the
primitive estimate instead. Fixed factorial shifts and the uniform
comparison for $\beta_g$ turn these bounds into
\eqref{eq:s08:marked-bound}. If $h=0$, use
Lemma~\ref{lmm:s06:vertex-bound} at zero cotangent parameters.
The finitely many stable low genera, including $(g,h)=(1,1)$,
are absorbed by enlarging the constants. No bound for a growing
number of markings is asserted here.
\end{prf}
